\documentclass[12pt]{article}
\usepackage[utf8]{inputenc}
\usepackage[T1]{fontenc}
\usepackage{hyperref}
\title{Kulina Pano Language Resources for K-12 Curriculum Development}
\author{Compiled by Claude Sonnet 5.5 (claude-sonnet-5-5) for Andr\'as Kornai}
\date{}
\begin{document}
\maketitle

\section{Introduction}
Kulina Pano is an obsolescent Panoan language of the Javari River valley in Brazil. Wikipedia reports 32 native speakers (2007), living in the village of Aldeia Pedro Lopes on the middle Curu\c{c}\'a River, Amazonas. Joshua Project gives a much larger population figure and places the people mostly in Acre. These sources disagree. Another Panoan language, Olivença Kulina, is also called Kulina and is extinct, so the two should not be confused. My searches found no native-language teaching materials.

\section{Primary Educational Resources}
I found no primers, textbooks or school programmes for Kulina Pano.

\section{Linguistic and Reference Materials}
\begin{itemize}
\item Wikipedia, Kulina language (Panoan): \url{https://en.wikipedia.org/wiki/Kulina_language_(Panoan)}
\item Wikipedia, Culina language (a different language, for disambiguation): \url{https://en.wikipedia.org/wiki/Culina_language}
\item Wikipedia, Panoan languages: \url{https://wikipedia.com/wiki/Panoan}
\item Numerals page for Kulina Pano, Max Planck Institute for Evolutionary Anthropology: \url{https://lingweb.eva.mpg.de/channumerals/Kulina-Pano.htm}
\end{itemize}

\section{Religious and Community Resources}
Joshua Project has a people-group profile, which is a mission-oriented source whose figures differ from the others: \url{https://joshuaproject.net/people_groups/21571}. I found no scripture translations and no indigenous community organisation.

\section{Digital and Audio Resources}
I found no audio, video or software.

\section{Curriculum Development Context}
With about 32 speakers, mostly older adults, a curriculum would need documentation work with the community first. School-based use is not known. Contact with Brazilian indigenous education bodies and with Panoan specialists would be the next step. Searches of Brazilian sources were not done here.

\end{document}